\documentclass{tufte-handout}

%\geometry{showframe}% for debugging purposes -- displays the margins

\usepackage[spanish, es-tabla]{babel}
\usepackage{amsmath}

% Set up the images/graphics package
\usepackage{gensymb}
\usepackage{physics}

\usepackage{graphicx}
\setkeys{Gin}{width=\linewidth,totalheight=\textheight,keepaspectratio}
\graphicspath{{graphics/}}

\title[Prácticas Química Física II: Normativa]{
Laboratorio Química Física II: Normativa}
%\author{David De Sancho}
\date{}  % if the \date{} command is left out, the current date will be used

% The following package makes prettier tables.  We're all about the bling!
\usepackage{booktabs}

% The units package provides nice, non-stacked fractions and better spacing
% for units.
\usepackage{units}

% The fancyvrb package lets us customize the formatting of verbatim
% environments.  We use a slightly smaller font.
\usepackage{fancyvrb}
\fvset{fontsize=\normalsize}

% Small sections of multiple columns
\usepackage{multicol}

% Provides paragraphs of dummy text
\usepackage{lipsum}

% These commands are used to pretty-print LaTeX commands
\newcommand{\doccmd}[1]{\texttt{\textbackslash#1}}% command name -- adds backslash automatically
\newcommand{\docopt}[1]{\ensuremath{\langle}\textrm{\textit{#1}}\ensuremath{\rangle}}% optional command argument
\newcommand{\docarg}[1]{\textrm{\textit{#1}}}% (required) command argument
\newenvironment{docspec}{\begin{quote}\noindent}{\end{quote}}% command specification environment
\newcommand{\docenv}[1]{\textsf{#1}}% environment name
\newcommand{\docpkg}[1]{\texttt{#1}}% package name
\newcommand{\doccls}[1]{\texttt{#1}}% document class name
\newcommand{\docclsopt}[1]{\texttt{#1}}% document class option name

\begin{document}

\maketitle% this prints the handout title, author, and date

%\begin{abstract}
%\noindent En esta primera práctica vamos a estudiar el espectro de rotación-vibración
%de las moléculas presentes en el aire que respiramos. Para ello usaremos
%espectroscopía en el infrarrojo, que permite estudiar la interacción de la
%materia con la \textbf{radiación infrarroja} ($\lambda=0.7$ $\mu$m$-1$ mm) 
%y así obtener espectros vibracionales con estructura fina rotacional.
%\end{abstract}

%\printclassoptions
\section{Normas en el Laboratorio}
En el laboratorio existe un material común, que incluye balanzas, secadores de vacío,
estufa, vitrinas, secadores de aire caliente, contenedores de agua, calefactores,
agitadores, ordenadores, que deben utilizarse de forma ordenada y limpia, procurando
una máxima utilización para todo el personal. En el caso de reactivos se pondrá un particular
interés en no contaminarlos mediante el uso de pipetas o espátulas sucias.

Todas las disoluciones, una vez utilizadas, deberán ser obligatoriamente vertidas en el
contenedor de residuos del laboratorio. A continuación, se limpiará el material
de manera convencional con jabón. El uso de la acetona como líquido de secado
deberá utilizarse sólo después de eliminado el jabón.

El laboratorio dispone de tres aparatos, dos espectrofotómetros UV-VIS y un espectrofotómetro FTIR,
cuya manipulación inicial requiere la presencia obligada del profesor responsable para su puesta en
marcha, evitando cualquier uso por parte del alumno antes de que
dicha puesta a punto se realice. El material correspondiente a cada aparato deberá ser
limpiado después del uso y se dejará en el mismo puesto.

\section{Seguridad en el Laboratorio}
Además de las normas básicas de seguridad a las que el alumno deberá estar familiarizado,
como el uso obligatorio de bata, gafas, guantes, se requiere poner énfasis en la
manipulación de los materiales típicos cuando se trabaja con disoluciones que contienen
disolventes orgánicos.
La utilización de disolventes orgánicos volátiles requiere que su manipulación se lleve a cabo
en vitrina o, en su defecto, en el puesto de trabajo, evitando que los recipientes que los
contienen estén abiertos durante mucho tiempo. Asimismo, bajo ningún concepto se debe
aspirar con la boca las pipetas.

La realización de disoluciones, que generalmente requieren agitación y calentamiento, deberá
hacerse en la vitrina. En caso de calentamiento, deberá evitarse que el líquido hierva, ya que
algunos son inflamables, se realizará sin tapón y deberá vigilarse constantemente por el
alumno.
La manipulación de ácidos se hará obligatoriamente en la vitrina.

Debido a la continua utilización de líquidos orgánicos inflamables, queda absolutamente
prohibida la introducción de material de vidrio con trazas de los mismos en la estufa de
secado, en la que sólo se introducirá material no graduado con trazas de agua.

Una última recomendación es que el alumno no dude en consultar con el profesor responsable cualquier duda o problema que se plantee.

\end{document}